证明 2.5. 作为 $k[[t]]$-代数, 局部环 $\mathcal O_{X,x}$ 的完备化同构于商
\[
\widehat{\mathcal O}_{X,x}
=k[[t,x_1,\ldots,x_{n+1}]]/(f).
\]
隐函数定理 XV, 1.1.2（即 [1], 引理 5.10）说明, 存在整数 $N$, 使得只要有 $k[[t]]$-代数同构
\[
\widehat{\mathcal O}_{Y,y_0}
\simeq k[[t,x_1,\ldots,x_{n+1}]]/(f_1)
\]
且 $f_1-f$ 属于 $k[[t,x_1,\ldots,x_{n+1}]]$ 的极大理想的 $N$ 次幂 $\mathfrak m^N$, 条件 2.5(b) 就成立.

设 $(t_0,t_1)$ 和 $(x_0,\ldots,x_{n+1})$ 分别为 $\mathbb P^1$ 与 $\mathbb P^{n+1}$ 上通常的齐次坐标系. 取 $Y/\mathbb P^1$ 为 $\mathbb P^1\times\mathbb P^{n+1}$ 中由方程
\begin{equation}\tag{2.5.1}
g(t_0,t_1;x_0,\ldots,x_{n+1})=0,
\end{equation}
定义的子簇, 其中该式关于 $(t_0,t_1)$ 为 $d'$ 次齐次, 关于各 $x_i$ 为 $d''$ 次齐次. 更内蕴地说, $g$ 是
\[
\mathcal O(d',d'')=\operatorname{pr}_1^*\mathcal O(d')
\otimes\operatorname{pr}_2^*\mathcal O(d'').
\]
的一个截面. 若
\begin{equation}\tag{2.5.2}
g(1,t;1,x_1,\ldots,x_{n+1})\equiv f\pmod{\mathfrak m^N}.
\end{equation}
则条件 (b) 成立.

固定 $d'$ 和 $d''$. 形如 (2.5.1) 的截面 $g$ 是某个仿射空间 $A$ 的 $k$-有理点. 对 $u\in\mathbb P^1\times\mathbb P^{n+1}$, 以 $B(u,r)$ 表示以 $\mathcal O(d',d'')$ 的截面在 $u$ 处的 $r$ 阶喷射为点的仿射空间. 置 $u_0=(1,0;1,0,\ldots,0)$. 若 $u\ne u_0$, 则当 $d'$ 和 $d''$ 充分大时, 始终为线性的限制映射
\begin{equation}\tag{2.5.3}
A\longrightarrow B(u_0,N)\times B(u,2)
\end{equation}
是满射.

设 $C(u)\subset B(u,2)$ 为如下 $2$ 阶喷射的集合: 它们或者是零芽的 $2$ 阶喷射, 或者是满足下述条件的芽 $\widetilde h$ 的二阶喷射:
\begin{enumerate}
\item[(a)] $u$ 属于由 $\widetilde h=0$ 定义的簇 $V$;
\item[(b)] $u$ 既不是 $V/\mathbb P^1$ 中通过 $u$ 的纤维的光滑点, 也不是其普通二次点.
\end{enumerate}
容易验证, $C(u)$ 在 $B(u,2)$ 中为余维 $n+3$ 的闭集.

设 $D$ 为 $A$ 中由满足 (2.5.2) 的 $g$ 所成的仿射子空间. 对 $u\ne u_0$, 设 $D(u)$ 为 $D$ 中由零截面 $g$ 以及满足下述条件的 $g$ 所成的子空间: $g(u)=0$, 且 $u$ 既不是由 $g=0$ 定义的簇 $V$ 中通过 $u$ 的纤维的光滑点, 也不是其普通二次点. 若 $d'$ 和 $d''$ 充分大, 使 (2.5.3) 满射, 则 $D(u)$ 在 $D$ 中的余维为 $n+3$. 因此所有 $u\ne u_0$ 的 $D(u)$ 之并的余维至少为
\[
1=(n+3)-(n+2),
\]
而 $D$ 的一般元素 $g$ 不属于任何 $D(u)$. 它就解决了问题.

\subsection*{参考文献}

\begin{enumerate}[label={[\arabic*]},leftmargin=3em]
\item M. Artin, ``Algebraic approximation of structures over complete
local rings,'' \emph{Publ. Math. IHES} \underline{36} (1969)
（Zariski 纪念卷）.
\item L\^e D\~ung Tr\'ang, \emph{Singularit\'es isol\'ees des
hypersurfaces complexes}, Publ. Centre Math. \'Ecole Polytechnique,
1969 年 5 月.
\item J. Milnor, \emph{Singular Points of Complex Hypersurfaces},
Princeton University Press, 1968.
\item M. Raynaud, ``Caract\'eristique d'Euler--Poincar\'e d'un
faisceau et cohomologie des vari\'et\'es ab\'eliennes,'' S\'eminaire
Bourbaki 286, 1965 年 2 月.
\item J.-P. Serre, \emph{Corps locaux}, Publ. Inst. Math. Univ.
Nancago VIII, Hermann, 1962.
\item J.-P. Serre, ``SGA 5, IX. Introduction \`a la th\'eorie de
Brauer,'' 收录于 \emph{Repr\'esentations lin\'eaires des
groupes}, 第 3 部分, Hermann, Paris, 1967.
\end{enumerate}
